% Paste this section into an existing document after loading amsmath and amssymb.
% The companion references.tex supplies the citation keys used below.
\section{Ionisation balance in the solar atmosphere}
\label{sec:ionisation_physics}

The ionisation state of an element determines the fraction of its population
that can respond to the ponderomotive force. For FIP fractionation, we therefore
require the relative populations of neutral atoms and ions throughout the
chromosphere. These populations depend on electron collisions, the radiation
field and charge exchange with the background plasma.

We describe below a rate-based treatment on a prescribed atmosphere. The
equations specify the physical processes required for this treatment; their
inclusion in a particular FIPpy calculation depends on the selected backend
and atomic data. In the existing metal prescription, FIPpy combines Saha and
CHIANTI equilibrium fractions. The rate equations developed here describe the
proposed replacement.

\subsection{Population fractions and rate coefficients}
\label{subsec:population_definitions}

For an element \(X\) with atomic number \(Z_X\), let \(n_{X,q}\) denote the
number density in charge state \(q\), where \(q=0\) denotes the neutral atom.
The total element density and charge-state fractions are
\begin{equation}
 n_X=\sum_{q=0}^{Z_X}n_{X,q},
 \qquad
 f_{X,q}=\frac{n_{X,q}}{n_X},
 \qquad
 \sum_{q=0}^{Z_X}f_{X,q}=1.
 \label{eq:ion_population_definitions}
\end{equation}
The total ionised fraction is then
\begin{equation}
 \xi_X=\sum_{q=1}^{Z_X}f_{X,q}=1-f_{X,0}.
 \label{eq:ion_total_fraction}
\end{equation}
Here \(n_{X,0}\) includes all populated bound levels of the neutral atom.
If excited neutral levels are present, one minus the ground-state fraction
would overestimate \(\xi_X\).

We distinguish a rate coefficient from a transition rate per particle.
For a binary electron collision, the coefficient \(C_{X,q}\) has units of
\(\mathrm{cm^3\,s^{-1}}\), while the rate experienced by one target particle is
\begin{equation}
 P_{X,q}^{\mathrm{e,ion}}=n_{\mathrm e}C_{X,q},
 \qquad
 [P_{X,q}^{\mathrm{e,ion}}]=\mathrm{s^{-1}}.
 \label{eq:ion_coefficient_rate}
\end{equation}
The number of such transitions per unit volume and time is
\(n_{X,q}n_{\mathrm e}C_{X,q}\).
Photoionisation is described directly by a rate per target particle,
\(\Gamma_{X,q}\), with units of \(\mathrm{s^{-1}}\).

\subsection{Electron-impact ionisation}
\label{subsec:electron_ionisation}

Direct ionisation occurs when a free electron transfers enough energy to
remove a bound electron:
\begin{equation}
 X^{q+}+\mathrm e^-
 \longrightarrow X^{(q+1)+}+2\mathrm e^-.
 \label{eq:ion_direct_reaction}
\end{equation}
The rate coefficient depends on the collision cross-section and the electron
energy distribution. The standard temperature-dependent coefficients used here
assume a Maxwellian distribution. At low electron temperatures, only a small
fraction of the electrons have energies above the relevant ionisation threshold.

An electron collision can also excite an ion into an unstable state above an
ionisation threshold. This state may decay by ejecting an electron:
\begin{align}
 X^{q+}+\mathrm e^-
 &\longrightarrow (X^{q+})^{**}+\mathrm e^-,
 \nonumber\\
 (X^{q+})^{**}
 &\longrightarrow X^{(q+1)+}+\mathrm e^-.
 \label{eq:ion_ea_reaction}
\end{align}
This process is called excitation--autoionisation. Ordinary excitation to a
bound level changes the internal energy of the atom without changing its
charge state. Its subsequent transitions can nevertheless affect the
ionisation balance in a multilevel calculation.

The total electron-impact ionisation coefficient is
\begin{equation}
 C_{X,q}(T_{\mathrm e})
 =C_{X,q}^{\mathrm{DI}}(T_{\mathrm e})
 +C_{X,q}^{\mathrm{EA}}(T_{\mathrm e}),
 \label{eq:ion_total_electron}
\end{equation}
where DI and EA denote direct ionisation and excitation--autoionisation.
These contributions are available through the corresponding
\texttt{fiasco} properties \texttt{direct\_ionization\_rate} and
\texttt{excitation\_autoionization\_rate}. Their combined coefficient is
\texttt{ionization\_rate} \cite{fiasco_rates}.

\subsection{Radiative and dielectronic recombination}
\label{subsec:electron_recombination}

Radiative recombination occurs when an ion captures a free electron and emits
a photon:
\begin{equation}
 X^{(q+1)+}+\mathrm e^-
 \longrightarrow X^{q+}+h\nu.
 \label{eq:ion_rr_reaction}
\end{equation}
The emitted photon carries away the energy required for the electron to remain
bound. The coefficient \(\alpha_{X,q+1}^{\mathrm{RR}}(T_{\mathrm e})\)
describes capture from the recombining charge state \(q+1\). In the usual
spontaneous-recombination treatment, this coefficient depends on the electron
temperature and atomic structure. An incident radiation field is not required
for the capture to occur.

Dielectronic recombination involves capture of a free electron together with
excitation of an electron already bound to the ion:
\begin{align}
 X^{(q+1)+}+\mathrm e^-
 &\longrightarrow (X^{q+})^{**},
 \nonumber\\
 (X^{q+})^{**}
 &\longrightarrow X^{q+}+h\nu.
 \label{eq:ion_dr_reaction}
\end{align}
The intermediate state can either autoionise, returning an electron to the
continuum, or emit a stabilising photon. Recombination succeeds when the ion
stabilises into a bound state. The term dielectronic refers to the captured
electron and the excited bound electron.

When the dataset provides separate radiative and dielectronic contributions,
the total coefficient is
\begin{equation}
 \alpha_{X,q+1}
 =\alpha_{X,q+1}^{\mathrm{RR}}
 +\alpha_{X,q+1}^{\mathrm{DR}}.
 \label{eq:ion_total_recombination}
\end{equation}
In \texttt{fiasco}, we use \texttt{recombination\_rate} for the preferred total.
Some ions have a dedicated total-recombination fit, for which the returned
coefficient can differ from a manual sum of the two component properties
\cite{fiasco_rates}.

At finite electron densities, collisions can reionise highly excited states
formed during dielectronic recombination before they complete their radiative
decay. This reduces the effective recombination coefficient. Density suppression
must be applied consistently with the adopted data, since a density-dependent
effective coefficient may already include this correction
\cite{nikolic2018}.

\subsection{Photoionisation and the radiation field}
\label{subsec:photoionisation}

Photoionisation occurs when an atom or ion absorbs a photon with sufficient
energy to remove an electron:
\begin{equation}
 X^{q+}+h\nu
 \longrightarrow X^{(q+1)+}+\mathrm e^-,
 \qquad h\nu\geq\chi_{X,q}.
 \label{eq:ion_photo_reaction}
\end{equation}
Here \(\chi_{X,q}\) is the ionisation energy of the initial state.
For a ground-state calculation it is the ground-state ionisation threshold;
excited levels have their own thresholds and cross-sections.

The local radiation can contain photons emitted in the photosphere,
transition region and corona, as well as emission within the chromosphere.
These photons allow ionisation to occur in a cool parcel whose thermal
electrons rarely reach the required energy. The relative importance of each
source depends on wavelength, geometry and attenuation. Photospheric radiation
therefore need not dominate the ionisation of every element at every height
\cite{laming2015,dufresne2021}.

Let \(I_\nu(s,\Omega)\) denote the specific intensity at position \(s\), in
direction \(\Omega\). The mean intensity is
\begin{equation}
 J_\nu(s)=\frac{1}{4\pi}\int I_\nu(s,\Omega)\,\mathrm d\Omega.
 \label{eq:ion_mean_intensity}
\end{equation}
The photoionisation rate per target particle is
\begin{equation}
 \Gamma_{X,q}(s)
 =4\pi\int_{\nu_{\mathrm{th},X,q}}^\infty
 \frac{J_\nu(s)}{h\nu}\,
 \sigma_{X,q}(\nu)\,\mathrm d\nu,
 \qquad
 h\nu_{\mathrm{th},X,q}=\chi_{X,q}.
 \label{eq:ion_photo_frequency}
\end{equation}
Dividing the intensity by \(h\nu\) converts energy intensity into photon-number
intensity. Multiplication by the photoionisation cross-section
\(\sigma_{X,q}\) gives the absorption rate for one target particle.
The result already has units of \(\mathrm{s^{-1}}\), so no factor of
\(n_{\mathrm e}\) is required.

For a spectrum sampled in wavelength, the equivalent expression is
\begin{equation}
 \Gamma_{X,q}(s)
 =4\pi\int_0^{\lambda_{\mathrm{th},X,q}}
 J_\lambda(s)\frac{\lambda}{hc}
 \sigma_{X,q}(\lambda)\,\mathrm d\lambda,
 \qquad
 \lambda_{\mathrm{th},X,q}=\frac{hc}{\chi_{X,q}}.
 \label{eq:ion_photo_wavelength}
\end{equation}
We use spectral densities consistently, with
\(J_\lambda=J_\nu c/\lambda^2\) when both are expressed per corresponding
physical units of wavelength and frequency. For example,
\(J_\lambda\) in
\(\mathrm{erg\,s^{-1}\,cm^{-2}\,sr^{-1}\,nm^{-1}}\),
\(\sigma\) in \(\mathrm{cm^2}\), and integration over nm give a rate in
\(\mathrm{s^{-1}}\).

Analytical ground-state cross-sections can be obtained from
\cite{verner1996}. The photon cross-section in
Eq.~\eqref{eq:ion_photo_frequency} is distinct from the electron-impact
cross-section used to evaluate \(C_{X,q}\). A numerical calculation must
resolve the ionisation threshold and the spectral features that contribute
to the integral.

\subsection{Opacity, attenuation and local emission}
\label{subsec:ion_radiative_transfer}

Opacity describes how radiation changes as it propagates through the
atmosphere. Along a ray with path coordinate \(\ell\), the transfer equation is
\begin{equation}
 \frac{\mathrm d I_\nu}{\mathrm d\ell}
 =-\kappa_\nu I_\nu+j_\nu,
 \label{eq:ion_transfer}
\end{equation}
where \(\kappa_\nu\) is the extinction coefficient in
\(\mathrm{cm^{-1}}\) and \(j_\nu\) is the emission coefficient per unit volume
and solid angle. Scattering introduces contributions to both extinction and
emission. Their treatment depends on the adopted transfer model.

If we neglect emission and scattering along the ray, its intensity reduces to
\begin{equation}
 I_\nu(\ell)=I_{\nu,0}\exp[-\tau_\nu(\ell)],
 \qquad
 \tau_\nu(\ell)=\int_{\ell_0}^{\ell}\kappa_\nu(\ell')\,\mathrm d\ell'.
 \label{eq:ion_attenuation}
\end{equation}
An H-only continuum approximation uses
\begin{equation}
 \kappa_\nu^{\mathrm{H\,I}}
 =n_{\mathrm{H^0}}\sigma_{\mathrm{H\,I}}(\nu).
 \label{eq:ion_h_opacity}
\end{equation}
This prescription requires both the incoming intensity and the neutral
hydrogen column. Opacity alone does not specify the number of photons present.
We first calculate the transmitted radiation, average over directions to
obtain \(J_\nu\), and then evaluate
Eq.~\eqref{eq:ion_photo_frequency}.

The ground-state H~I bound-free cross-section vanishes below its ionisation
threshold. Metal continua and spectral lines can still provide opacity at
these photon energies. In particular, H~I continuum attenuation does not
describe resonant scattering in Ly\(\alpha\). Locally trapped Ly\(\alpha\)
requires an emission and escape prescription or a transfer calculation.
Its ability to photoionise a species also depends on the threshold of the
initial atomic level \cite{laming2015,dufresne2025}.

An expression such as \(J_\nu=J_{\nu,0}\exp(-\tau_\nu)\) represents an
effective-ray approximation when used directly for the mean intensity.
A directional calculation instead attenuates each incoming ray before
performing the angular average. The ray path through the atmosphere must be
specified separately from the coordinate along a curved magnetic loop.

\subsection{Irradiance and boundary illumination}
\label{subsec:ion_boundary_radiation}

An observed spectral irradiance \(F_{\lambda,d}\) at distance \(d\) is an energy
flux per unit wavelength. A conversion to local mean intensity requires
assumptions about the source. For a uniformly bright photospheric surface,
the outward surface flux and surface specific intensity satisfy
\begin{equation}
 F_{\lambda,\mathrm{surf}}
 =F_{\lambda,d}\left(\frac{d}{R_\odot}\right)^2,
 \qquad
 I_{\lambda,\mathrm{surf}}
 =\frac{F_{\lambda,\mathrm{surf}}}{\pi}.
 \label{eq:ion_irradiance_conversion}
\end{equation}
Immediately above the surface, the emitting hemisphere occupies
\(2\pi\) steradians. Its contribution to the mean intensity is therefore
\begin{equation}
 J_{\lambda,\mathrm{ph},0}
 =\frac{I_{\lambda,\mathrm{surf}}}{2}
 =\frac{F_{\lambda,d}}{2\pi}
 \left(\frac{d}{R_\odot}\right)^2.
 \label{eq:ion_photospheric_boundary}
\end{equation}
The factor \(1/2\) follows from the angular average of the specific intensity
over the hemisphere. It should not be applied to an irradiance without this
conversion.

Downward coronal illumination requires a separate geometrical prescription.
An irradiance measured outside the Sun does not uniquely determine the
radiation incident on an individual chromospheric footpoint.
We must specify the illuminating spectrum, its angular distribution and the
paths used for attenuation. The wavelength interval must cover all photons
that contribute to the selected ionisation thresholds, without gaps between
separately assembled spectral components.

\subsection{Charge exchange with hydrogen}
\label{subsec:ion_charge_exchange}

Charge exchange transfers a bound electron during an atom--ion collision.
For a neutral element and a proton, the ionising direction is
\begin{equation}
 X^0+\mathrm H^+
 \longrightarrow X^++\mathrm H^0.
 \label{eq:ion_cx_forward}
\end{equation}
The proton captures an electron from \(X\), leaving the element ionised and
hydrogen neutral. The reverse reaction is
\begin{equation}
 X^++\mathrm H^0
 \longrightarrow X^0+\mathrm H^+.
 \label{eq:ion_cx_reverse}
\end{equation}
Both reactions conserve electric charge. Their rates depend on the densities
of the appropriate collision partners:
\begin{equation}
 P_X^{\mathrm{cx,ion}}
 =n_{\mathrm{H^+}}k_X^{\mathrm{cx,ion}},
 \qquad
 P_X^{\mathrm{cx,rec}}
 =n_{\mathrm{H^0}}k_X^{\mathrm{cx,rec}}.
 \label{eq:ion_cx_rates}
\end{equation}
The coefficients \(k_X\) have units of \(\mathrm{cm^3\,s^{-1}}\).
Their temperature dependence describes the relative motion of the heavy
collision partners. The electron temperature is an appropriate substitute
only under the corresponding thermal assumptions.

Oxygen is a useful example because its first ionisation energy is close to
that of hydrogen. Rapid O--H exchange can therefore couple the oxygen
ionisation balance closely to the hydrogen populations
\cite{dufresne2021}. In the exchange-only, two-stage equilibrium limit,
\begin{equation}
 \frac{n_{\mathrm{O^+}}}{n_{\mathrm{O^0}}}
 =
 \frac{n_{\mathrm{H^+}}}{n_{\mathrm{H^0}}}
 \frac{k_{\mathrm O}^{\mathrm{cx,ion}}}
      {k_{\mathrm O}^{\mathrm{cx,rec}}}.
 \label{eq:ion_oxygen_exchange_limit}
\end{equation}
Exact equality of the H and O ionisation fractions requires further conditions
on the coefficient ratio and on the competing processes.

Charge exchange also affects other elements, for which reliable
species-dependent rates are needed. The revised sulphur treatment discussed by
\cite{laming2021} illustrates the sensitivity to the adopted data.
These reaction coefficients are separate from elastic momentum-transfer
collision frequencies used in the fractionation calculation.

\subsection{Combined two-stage statistical equilibrium}
\label{subsec:ion_two_stage}

For a first approximation, we retain only \(X^0\) and \(X^+\) and treat the
background atmosphere as prescribed. The total ionisation and neutralisation
rates per target particle are
\begin{align}
 R_{\mathrm{ion},X}
 &=n_{\mathrm e}C_{X,0}+\Gamma_{X,0}
   +n_{\mathrm{H^+}}k_X^{\mathrm{cx,ion}},
 \label{eq:ion_combined_forward}\\
 R_{\mathrm{neu},X}
 &=n_{\mathrm e}\alpha_{X,1}
   +n_{\mathrm{H^0}}k_X^{\mathrm{cx,rec}}.
 \label{eq:ion_combined_reverse}
\end{align}
All terms in these expressions have units of \(\mathrm{s^{-1}}\).
For this closed two-stage system, local statistical equilibrium requires
\begin{equation}
 n_{X,0}R_{\mathrm{ion},X}
 =n_{X,1}R_{\mathrm{neu},X}.
 \label{eq:ion_balance_two}
\end{equation}
Using \(n_X=n_{X,0}+n_{X,1}\), we obtain
\begin{equation}
 \xi_X^{\mathrm{SE}}
 =\frac{R_{\mathrm{ion},X}}
 {R_{\mathrm{ion},X}+R_{\mathrm{neu},X}}.
 \label{eq:ion_fraction_two}
\end{equation}
We use the term statistical equilibrium because the balance includes
collisional and radiative processes as well as charge exchange.
Photoionisation can be dominant or subdominant within the same formulation.

At a fixed position, the relative ionising contributions can be assessed by
comparing \(\Gamma_{X,0}\), \(n_{\mathrm e}C_{X,0}\) and
\(n_{\mathrm{H^+}}k_X^{\mathrm{cx,ion}}\). A large contribution to a local
rate does not by itself establish a large effect on the integrated FIP bias;
the neutral fraction must change where fractionation occurs.

\subsection{CHIANTI coefficients through \texttt{fiasco}}
\label{subsec:ion_fiasco}

For a neutral-to-singly-ionised calculation, the electron-impact coefficient
must describe ionisation out of \(X^0\), whereas the recombination coefficient
must describe capture by \(X^+\). Thus, oxygen requires
\texttt{ionization\_rate} from \texttt{O I} and
\texttt{recombination\_rate} from \texttt{O II}.
For a plotted \texttt{Fe XVI} object, the ionisation coefficient describes
Fe~XVI to Fe~XVII, while the recombination coefficient describes Fe~XVI to
Fe~XV. They describe losses from the same initial charge state into different
neighbouring stages \cite{fiasco_rates}.

The \texttt{fiasco} coefficients supply the electron-driven part of
Eqs.~\eqref{eq:ion_combined_forward} and \eqref{eq:ion_combined_reverse}.
The local photon rate requires \(J_\nu\) and photoionisation cross-sections;
charge exchange requires the corresponding reaction data and H populations.
The CHIANTI version and coefficient validity ranges must be recorded.
Support for a database version does not establish that a given coefficient
property includes every density or metastable correction available in that
database.

\subsection{LTE, coronal equilibrium and time dependence}
\label{subsec:ion_equilibrium_limits}

Under local thermodynamic equilibrium, the relative charge populations follow
the Saha equation. For neutral and singly ionised stages,
\begin{equation}
 \frac{n_{X,1}}{n_{X,0}}
 =\frac{2}{n_{\mathrm e}}
 \left(\frac{2\pi m_{\mathrm e}k_{\mathrm B}T}{h^2}\right)^{3/2}
 \frac{U_{X,1}(T)}{U_{X,0}(T)}
 \exp\left(-\frac{\chi_{X,0}}{k_{\mathrm B}T}\right),
 \label{eq:ion_saha}
\end{equation}
where \(U_{X,q}\) is the partition function of charge state \(q\).
Replacing the partition functions by constant ground-state statistical weights
is an additional approximation. Its accuracy must be checked for the element
and temperature range of interest.

In the low-density coronal-equilibrium approximation, the balance is set by
electron-impact ionisation and recombination. In the two-stage limit,
Eq.~\eqref{eq:ion_fraction_two} becomes
\begin{equation}
 \xi_X^{\mathrm{CIE}}
 =\frac{C_{X,0}}{C_{X,0}+\alpha_{X,1}}.
 \label{eq:ion_cie_two}
\end{equation}
Its explicit density factors cancel when both coefficients depend on
temperature alone. A multistage coronal calculation instead solves for all
charge fractions and obtains \(\xi_X\) from Eq.~\eqref{eq:ion_total_fraction}.
Density-dependent coefficients can change this simple cancellation.

The rate model containing electron-impact ionisation, spontaneous radiative
and dielectronic recombination, photoionisation and charge exchange is still
an approximation. Recovering LTE requires consistent inverse processes and
detailed balance, including the relevant three-body and radiative terms.
Additional levels may be required for metastable and excited-state populations.

For fixed rates in a closed two-stage system, the time-dependent fraction obeys
\begin{equation}
 \frac{\mathrm d\xi_X}{\mathrm dt}
 =(1-\xi_X)R_{\mathrm{ion},X}
 -\xi_X R_{\mathrm{neu},X}.
 \label{eq:ion_time_two}
\end{equation}
Its relaxation time is
\begin{equation}
 \tau_X=\frac{1}{R_{\mathrm{ion},X}+R_{\mathrm{neu},X}}.
 \label{eq:ion_relaxation}
\end{equation}
Statistical equilibrium is appropriate when this relaxation is fast compared
with changes in the atmosphere and irradiation, and with transport through
the relevant gradients. Hydrogen can depart substantially from instantaneous
equilibrium in a dynamic chromosphere \cite{carlsson2002}.
A two-stage relaxation time does not describe all modes of a multilevel atom.

\subsection{Extension to multiple charge states}
\label{subsec:ion_multistage}

Let \(\mathcal I_q\) denote the total transition rate from charge \(q\) to
\(q+1\), and \(\mathcal R_q\) the rate from \(q\) to \(q-1\).
For a network containing adjacent-stage transitions, the local population
equations are
\begin{equation}
 \frac{\mathrm d f_{X,q}}{\mathrm dt}
 =\mathcal I_{q-1}f_{X,q-1}
  +\mathcal R_{q+1}f_{X,q+1}
  -(\mathcal I_q+\mathcal R_q)f_{X,q}.
 \label{eq:ion_multistage}
\end{equation}
Terms outside the allowed charge range are omitted, with
\(\mathcal R_0=\mathcal I_{Z_X}=0\).
In matrix form,
\begin{equation}
 \frac{\mathrm d\boldsymbol f_X}{\mathrm dt}
 =\boldsymbol M_X\boldsymbol f_X,
 \qquad
 \boldsymbol M_X\boldsymbol f_X=0
 \quad\hbox{in statistical equilibrium}.
 \label{eq:ion_rate_matrix}
\end{equation}
We impose the normalisation in Eq.~\eqref{eq:ion_population_definitions}.
Using a column-vector convention, the columns of \(\boldsymbol M_X\) sum to
zero. A level-resolved model enlarges this system to include bound-level
transitions within each charge state.

On a prescribed atmosphere, \(T_{\mathrm e}\), \(n_{\mathrm e}\), the H
populations and the radiation field are inputs. Such a calculation does not
automatically update the electron density or feed ionisation energy and
radiative losses back into the hydrodynamics. These couplings require further
equations and their own validation. The resulting populations can nevertheless
provide \(\xi_X\) for the FIP calculation, provided the adopted ion
momentum-transfer treatment remains appropriate for the populated charge states.
